\section{Data and Measurement} \label{sec:data_c}

To evaluate the impact of the financial education pilot, we use a combination of survey and administrative data to measure a broad set of outcomes over different time horizons. Survey data collected at the end of the 2015 school year capture short-term effects on students’ financial proficiency, hypothetical attitudes toward consumption and saving, and self-reported behaviors. The survey data were collected right after the program implementation finished, in line with other financial education programs targeting elementary and middle school students \citep{de2021effect, batty2020experiential, berry2018impact}. Administrative records provide short- and medium-term outcomes, including grade progression, standardized test scores, and dropout rates, observed up to three years after the intervention.


\subsection{Survey Data} \label{subsec:question}

The Center for Public Policy and Education Evaluation (CAEd), a center at the Federal University of Juiz de Fora specializing in standardized assessments and large-scale data collection, was responsible for implementing the survey, with students and teachers participating in the pilot. Data collection occurred during the final school days of the 2015 academic year, when CAEd distributed printed questionnaires directly in classrooms. Despite being conducted at the end of the school year, a period when some students may already be on vacation or have dropped out, student participation rates remained high, averaging approximately 80\%. Furthermore, as shown in \autoref{tab:studentscar}, there is no significant difference between participation rates of treatment and control groups across all grades.


Students completed three paper-based questionnaires during a two-hour session: (i) a financial literacy test, (ii) a questionnaire on hypothetical consumption and savings attitudes and self-reported behaviors, and (iii) a socioeconomic questionnaire. The content of each instrument was tailored to the students’ grade level and cognitive development. Fifth-, seventh-, and ninth-graders answered all three questionnaires in school. Third graders answered only the financial literacy test, which was read aloud by trained enumerators, given that students at this age are still developing reading and writing skills; the attitudes-and-behaviors questionnaire was not administered to this grade, and the socioeconomic questionnaire was sent to the students' homes for their legal guardians to fill out and return the following day.\footnote{Schools prepared the legal guardians by sending them a note before the test day. \autoref{tab:studentscar} shows that an average of 65\% of parents/guardians returned the questionnaire, with no statistically significant difference between treated and control students.}

The financial literacy test evaluated students’ understanding of concepts introduced in the financial education textbooks through multiple-choice questions aligned with grade-appropriate skills (\autoref{tab:descriptors}). The test includes items that assess students’ ability to interpret everyday financial information, such as utility bills or spending records, and to identify the purpose of documents like receipts, checks, and invoices. It also covers savings, expenses, consumption, waste, risk, return, financial planning, and investment. Additional questions assess students’ ability to estimate values for simple financial projects, distinguish between paid and unpaid work, interpret financial graphs and tables, and extract relevant financial information from media texts. Not all skills are assessed in all grades; the test content is adapted to the cognitive level of each group, with some skills appearing only in earlier or later grades. CAEd used student responses to construct a financial proficiency index using Item Response Theory (IRT), which accounts for the difficulty of each question, generates a standardized score comparable across grades, and allows us to assess the impact of the intervention in terms of standard deviations.

The second questionnaire measures students’  hypothetical attitudes toward consumption and saving and collects a set of self-reported behaviors. The attitudinal part consists of fifteen statements, seven on consumption and eight on saving, each printed with a four-point agreement scale running from “Totally agree” to “Totally disagree” (\autoref{tab:cs_index}). The statements are worded in both directions: agreeing with “I plan before spending my money” indicates a prudent attitude, while agreeing with “I see no problem in owing money” indicates the opposite. We reverse the scale of the positively worded statements so that in every item the highest value corresponds to the most prudent response. We summarize each block by its first principal component, standardized to mean zero and unit variance, and refer to the two resulting measures as the consumption and savings indices, an approach similar to \cite{Bruhn2016}. Since this second questionnaire was not administered in third-grade, the consumption and savings indices, like the self-reported behaviors below, are available only for grades 5, 7, and 9. These are non-incentivized, self-reported measures of what students say about money.

The second questionnaire also collects self-reported behaviors: whether the student keeps a piggy bank, receives an allowance, holds a bank account, a debit card or a credit card, and whether they discuss money with their parents or with their friends. We analyze these separately, combining only the bank account, debit card, and credit card items into a single indicator of access to financial services.

The last questionnaire asks about students' socioeconomic background, such as the educational attainment of their mothers, and whether they benefit from the national conditional cash transfer program, \textit{Bolsa Família}. As expected, \autoref{tab:studentscar} shows no significant differences between the socioeconomic characteristics of treated and control students.

The teachers' questionnaire investigates their wage level, experience, and the incorporation of financial education textbooks in the curriculum, among other questions. \autoref{tab:teacherscar} shows that there are no significant differences between teachers in treatment and control groups.

Overall, our analysis uses eight dependent variables from the survey data collected in the field. Three are standardized indices: (i) financial proficiency, computed by CAEd from the financial literacy test using Item Response Theory, and the (ii) consumption and (iii) savings attitude indices described above. The remaining five are self-reported behaviors: (iv) whether the student talks to their parents about financial subjects, (v) whether they talk to their friends about the same subjects, (vi) whether they keep a piggy bank, (vii) whether they have access to financial services, and (viii) whether they receive an allowance from their parents.


\subsection{Administrative Data}

Since 1995, every public and private K--12 school has reported to the annual School Census, implemented by the National Institute of Educational Studies and Research (INEP), a research agency under the Brazilian Ministry of Education. The Census records school infrastructure and students' outcomes such as age-grade distortion, grade promotion, retention, and dropout. Since 2005, INEP has also administered \textit{Prova Brasil}, a standardized assessment of reading and math proficiency given every two years to all students in the fifth and ninth grades of public schools.\footnote{Schools with at least 20 students enrolled in the evaluated grade. Students take the test at the end of the school year, between October and November.} Finally, since 1998 INEP has run the \textit{Exame Nacional do Ensino M\'edio} (\textit{ENEM}), a national exam taken at the end of high school. Unlike the other two, \textit{ENEM} is voluntary: students register on their own, and since 2009 most universities have used it for admission, which raised the number of registrations from about 157{,}000 in 1998 to 7.7 million in 2015.\footnote{INEP's \textit{ENEM} historical series records 157,221 registrations in the first edition (1998) and 7,746,057 confirmed registrations in 2015.}

Because of the limited budget and the richness of these public records, we did not collect a baseline student survey. Instead, we used the administrative data at the school level to check pre-treatment balance between treatment and control schools: \autoref{tab:pret_balancing} shows no significant differences in average proficiency in reading and math, grade promotion, dropout, retention, or age-grade distortion. Nonetheless, school-level balance has a limitation since not every student of a pilot school took part in the intervention: as explained in \autoref{sec:pilot}, principals selected the participating class at each grade, and treatment- and control-school principals could in principle have chosen classes on different, unobserved criteria.

To compare treated and control students rather than schools, in 2019 we asked INEP to merge the pilot roster with its student-level administrative records.\footnote{Because researchers cannot access identified student-level microdata, the merge itself was performed by INEP, using students' names, grades, and schools' identification numbers. We then ran all balance tests and regressions on the merged data ourselves inside INEP's secure facility (the cold room). \autoref{tab:balance_adm1} and \autoref{tab:balance_adm2} show no significant differences in the share of treated and control students that INEP located in its administrative datasets.} The merge serves two purposes. First, it lets us check balance at the student level. For the seventh- and ninth-grade cohorts, we recover the pre-pilot \textit{Prova Brasil} score (performance in reading and math) while in fifth grade. \autoref{tab:balance_adm2} shows that seventh-grade students are balanced on this pre-treatment score, whereas ninth-grade control students have significantly higher pre-treatment reading scores. If higher reading correlates with financial proficiency, this imbalance would bias our estimates downward. The third- and fifth-grade cohorts have no such baseline, because \textit{Prova Brasil} is first taken in the fifth grade.

Furthermore, the merge enriches the analysis with outcomes that exist only in the administrative records. \autoref{tab:inep_data} lists them by cohort: post-treatment proficiency in reading and math in \textit{Prova Brasil}, retention and dropout, observed annually from 2015 to 2018 in the School Census; and, for the ninth-grade cohort, high-school completion, \textit{\textit{ENEM}} participation, and \textit{ENEM} proficiency scores.\footnote{Because a pilot cohort is only tested in \textit{Prova Brasil} when it reaches a grade in which the exam is given, the timing of these outcomes differs across cohorts.}



\begin{table}[h!]
\centering
\begin{threeparttable}
\fontsize{9}{11}\selectfont
\caption{Administrative outcomes recovered for pilot students, by 2015 cohort}\label{tab:inep_data}
\begin{tabular*}{\textwidth}{@{\extracolsep{\fill}}l >{\centering\arraybackslash}p{2cm} >{\centering\arraybackslash}p{2cm} cccc@{}}
\toprule
 & Pre-treatment & \multicolumn{5}{c}{Post-treatment} \\
\cmidrule(lr){2-2}\cmidrule(lr){3-7}
 & Proficiency & Proficiency & Retention & Dropout & HS completion & \shortstack{Participation \&\\ proficiency} \\
Cohort & (\textit{Prova Brasil}) & (\textit{Prova Brasil}) & \multicolumn{3}{c}{(School Census)} & (\textit{\textit{ENEM}}) \\
\midrule
third graders in 2015 & --- & while in fifth grade in 2017 & 2015--2018 & 2015--2018 & --- & --- \\
fifth graders in 2015 & --- & while in fifth grade in 2015 & 2015--2018 & 2015--2018 & --- & --- \\
seventh graders in 2015 & while in fifth grade in 2013 & while in ninth grade in 2017 & 2015--2018 & 2015--2018 & --- & --- \\
ninth graders in 2015 & while in fifth grade in 2011 & while in ninth grade in 2015 & 2015--2018 & 2015--2018 & 2018 & 2018 \\
\bottomrule
\end{tabular*}
\begin{tablenotes}
\fontsize{9}{9}\selectfont
\item \textit{Notes}: Each row is a pilot cohort, defined by the grade the student attended in 2015, and each cell reports the year in which the outcome is observed for that cohort through the student-level merge with INEP's records. Dashes mark outcomes that do not exist for the cohort.
\item \textit{Prova Brasil} is administered only in the fifth and in the ninth grades, so each cohort is tested when it reaches one of these grades. The years in the table are those of students who progressed on schedule; a student who repeated a grade reaches the fifth or the ninth grade, and hence the test, in a later year. After the pilot, the fifth- and ninth-grade cohorts were tested in 2015, in the pilot year itself, and the third- and seventh-grade cohorts in 2017, when they reached the fifth and the ninth grade, respectively. Before the pilot, only the seventh- and ninth-grade cohorts had been tested, while in the fifth grade in 2013 and 2011, respectively. This pre-pilot score is the outcome we use to assess pre-pilot balance in students' performance between treated and control groups. The third- and fifth-grade cohorts had not reached the fifth grade before the pilot, so they have no baseline score.
\item \textit{Retention} and \textit{dropout} come from the annual School Census and are observed for every cohort in each year from 2015 to 2018. 
\item \textit{High-school completion} (School Census) and \textit{ENEM} participation and proficiency are observed in 2018 and exist only for the ninth-grade cohort, the only one to reach the end of high school within this window.
\item \textit{Source}: Authors' elaboration based on the pilot roster merged with INEP's identified student-level administrative records (School Census, \textit{Prova Brasil}, and \textit{ENEM}). INEP performed the merge; the authors ran all analyses inside INEP's secure facility.
\end{tablenotes}
\end{threeparttable}
\end{table}%
